% ===========================================================================
\section{Conclusion}
\label{sec:conclusion}
% ===========================================================================

Pakistan's waste sector emitted \SI{37.45}{\mega\tonne} \COe{} in FY2023--24 on AR5 metrics:
\SI{24.46}{\mega\tonne} from domestic wastewater, \SI{9.17}{\mega\tonne} from solid-waste
disposal and \SI{3.83}{\mega\tonne} from open burning. The sector is \SI{8.4}{\percent} of a
ten-sector national inventory and \SI{87}{\percent} methane, which makes it unusually
sensitive to the metric chosen and unusually relevant to near-term mitigation.

Two things generalise beyond this country.

\textbf{An activity figure shared between categories should be partitioned, not estimated
twice.} Solid-waste disposal and open burning describe two fates of one waste stream, and
computing them from one generation figure with $D + U = G$ asserted removes a class of
inconsistency that independent estimation cannot detect. The cost is nil and the gain is that
an argument about the collection rate becomes an argument about where emissions sit rather
than about how much waste exists. The same shape recurs wherever two categories divide one
national quantity.

\textbf{Where an assumption spans a factor approaching two and the evidence does not choose,
the defensible output is both numbers.} The sewer methane correction factor moves this sector
from \SI{37.45}{\mega\tonne} to \SI{28.38}{\mega\tonne}. The lower value happens to agree
better with EDGAR, and adopting it for that reason would have been a mistake: EDGAR is a model
resting on its own assumption about the same unmeasured quantity, so the agreement would have
recorded a shared default as though it were evidence about Pakistani drains. We report
\num{0.5} as central because the physical description of the drains is the stronger and more
specific evidence, publish \num{0.1} alongside, and name what would settle it.

Domestic wastewater remains a consequential uncertainty because sewer fate and discharge conditions are not observed comprehensively. Seasonal dissolved-methane and flux evidence could distinguish the retained structural cases. Open-burning pollutant factors and burn locations require separate evidence; a population-based waste-origin allocation is not an exposure estimate. The extended release advances the existing history rather than inventing annual management observations.
